% Generated from frozen plotting data; no fixed figure number for deployment.
\expandafter\def\csname fig4_revisedCaption\endcsname{Activation design and decoder-layer transform cost. (a) Activation representation budget versus WikiText-2 perplexity for Llama-3.2-3B; all 11 existing (g, m) configurations and the bf16 reference are retained. The budget is 4 + 16(m+1)/g activation bits/value, not E28 model-file size or whole-model effective bits. Point positions retain full precision; labels use two decimals. Filled circles/squares connect m=1 measurements; open circles/triangles retain additional directions. The horizontal bracket compares scale resolution at g=256 versus 128; the vertical bracket identifies the group-constant/null-space contribution at the shared 4.25-bit budget (Hadamard versus PrismQuant, g=128, m=1). PPL is averaged over three seeds and 64 chunks per seed; bf16 uses one reference. The 8B counterpart is in the appendix. (b) Activation-only PPL recovery is 100(PPL\_Hadamard\ensuremath{-}PPL\_k)/(PPL\_Hadamard\ensuremath{-}PPL\_bf16). All 15 measurements, including nonmonotone changes, remain visible. The k=8 operating point recovers 38.59\%, 35.53\%, and 59.12\% for Llama 3B, Llama 8B, and Qwen3-8B-Base. Rank labels are categories: qkv/down caps are respectively 24/64, 32/112, and 32/96, applied per site. Llama uses E11 (three seeds, 64 chunks); Qwen uses E18-v2 (one seed, 146 chunks), with bf16 weights/KV and group-128 activation quantization. The retained summary PPLs do not supply a common uncertainty estimate across these protocols; no interval or monotonic benefit is claimed. The green band marks only the k=8 category; its width is a visual operating-point highlight, not an uncertainty interval. (c) The original rank-cost measurements retain k=8 and k=32 for both Llama models and both Hadamard references (dashed), on NVIDIA RTX PRO 6000 Blackwell Server Edition at T=2048. Colors match the 3B/8B curves in b. The ordinate is 100 t\_transform/(t\_decoder\_layer+t\_transform), combining separately benchmarked durations; it is not measured whole-model decode overhead. No session uncertainty is available for this original benchmark. A larger standalone view with both exact Hadamard values is included in the appendix. Whole-model deployment evidence is presented separately. The accuracy protocols and random-weight E28 runtime protocol do not form a checkpoint-level joint accuracy--speed Pareto point.}

\expandafter\def\csname fig_deployment_efficiencyCaption\endcsname{Deployment efficiency of PrismQuant. The implementation strip shows compact R4, U = X A and Z = X H\_block \ensuremath{-} U B, with offline factors A/B and a direct X-to-B branch; box widths are not timing proportions. Kernel B outputs FP16 to the shared per-token symmetric INT4 quantizer and CUTLASS INT4 GEMM. (a) Eager prefill input throughput, normalized to the matched FP16 model/batch in the original E28-v2 uniform0..255 cohort, with 2048 input tokens per sequence. Batch size 1 or 16 means one or sixteen sequences processed together in one prefill call. The dashed line is FP16 = 1.00; source tables include absolute input tokens/s. (b) Matched private current-stream CUDA Graph decode latency for FP16, Hadamard, and PrismQuant. PrismQuant is 1.22\ensuremath{\times} faster than FP16 on 8B; 3B remains slower than FP16. (c) Peak allocated inference bytes divided by 1e9, from the same Graph mode, with 49.72\%/56.34\% PrismQuant savings relative to FP16 for 3B/8B. This is inference memory, not reserved/NVML memory or weight-file size; centers are the maximum of three independent-session peaks. (d) Candidate/reference wall-time ratios for prebound versus generic launch at T=1 (dispatch-sensitive) and the complete R4 path with TC-B versus shuffle-B at T=2048 (bulk execution). The labeled linear x-axis spans 0.25--1.05 and contains every session and reference point. The percentage label is the elapsed-time reduction, 100(1\ensuremath{-}r), computed from the unrounded central ratio r. Each comparison has its own gray reference at 1.00; this is not FP16 or Hadamard. Centers are medians of three paired-session ratios, not pooled timing ratios. These are local implementation ablations, not B-only or complete-model speedups. All T=1/2048/32768 cases, including unfavorable TC-B cases, are retained in the appendix. Deployment timing centers use pooled medians over 150 formal runs (three sessions, 50 runs each after 10 warmups); prefill ratios use the corresponding pooled throughput medians. Whiskers show min--max of the three session medians, paired within each session for ratios; memory whiskers show session-peak ranges. They are not 95\% confidence intervals, and tiny ranges are not enlarged. All performance models have random weights. Graph decode uses batch 1, prefix 2048, 128 causal steps with the first 8 discarded, and one preallocated 2176-token page. Independent page-64 correctness tests verify page crossing; the formal timing does not measure dynamic page allocation. FP16 is the performance reference, distinct from the accuracy panels' bf16 reference. The E28 per-token symmetric quantizer and KV interfaces are not the paper-native group-128 asymmetric full deployment. The frozen report discloses large-accumulator FP16-conversion stress-test failures in the existing backend; 6/6 private Graph correctness passes do not imply every numerical test passes.}

\expandafter\def\csname appendixCaption\endcsname{Appendix evidence. The original E17 v3 rank-cost figure retains four PrismQuant points (k=8 and k=32 for both models) and both k=8 Hadamard references, on NVIDIA RTX PRO 6000 Blackwell Server Edition, T=2048. The source statistic is 100 times the separately benchmarked transform duration divided by decoder-layer duration plus transform duration; it is not measured end-to-end decode overhead. The original Triton do\_bench protocol and selected configurations are preserved in E17V3\_DONE.json and the original timing CSVs; no session uncertainty is fabricated. The matched eager/Graph figure uses only the private current-stream panel, with 150 samples per method/mode, and shows each method's corresponding latency. The all-shapes figure retains all 30 kernel comparisons: five reference families, two models, and three token counts. Centers are medians of three paired-session wall-time ratios and whiskers show their min--max; gray marks denote each family's own reference at 1.00. The B comparison measures full R4 paths with different B implementations. E17-native rows here are separate local A40 microbenchmarks of asymmetric packed-output transforms and must not be confused with the original RTX PRO 6000 rank-cost figure or the E28 FP16-output interface. The 8B metadata plot uses the same E20 definition as Figure 4a and retains all 11 observations plus bf16.}
